% !TEX root = ../main.tex
\documentclass[../Main.tex]{subfiles}

\begin{document}

\chapter{Kom i gang med Python}
\label{app:python}

\intro{Python er et programmeringsspråk som er svært brukt i vitenskap, ingeniørfag og dataanalyse. Her gir vi en helt grunnleggende innføring i det du trenger for å kjøre kodeeksemplene i dette kompendiet. Du trenger ingen forkunnskaper.}

\section{Jupyter Notebook}

Den enkleste måten å skrive og kjøre Python på er via \textbf{Jupyter Notebook}. En notebook er et dokument der du kan blande kode, tekst og figurer. Slik kommer du i gang:

\begin{enumerate}
    \item Installer \textbf{Anaconda} fra \url{https://www.anaconda.com} (gratis). Dette installerer Python og Jupyter automatisk.
    \item Åpne \textbf{Anaconda Navigator} og klikk på \textit{Launch} under Jupyter Notebook.
    \item En nettleser åpnes. Naviger til mappen \texttt{TRES0312 Python} og åpne en av \texttt{.ipynb}-filene (samme filer ligger også på GitHub: \pyfolderlink).
    \item Klikk på en kodecelle og trykk \textbf{Shift+Enter} for å kjøre den.
\end{enumerate}

\begin{mymerk*}{Alle eksempler er klare til bruk}
Alle kodeeksemplene fra dette kompendiet finnes som ferdige Jupyter Notebook-filer i mappen \pyfolderlink. Du kan åpne dem direkte og kjøre dem uten å skrive noe selv — og deretter modifisere dem for å eksperimentere.
\end{mymerk*}

\section{Variabler og regneoperasjoner}

I Python lagrer vi verdier i \emph{variabler}. Du oppretter en variabel ved å skrive \texttt{navn = verdi}:

\begin{pythonkode}
m = 2.0    # masse i kg
v = 5.0    # fart i m/s
\end{pythonkode}
\pyfilenote{Starten av samme kodecelle som i \href{\pyrepobase/Kap2_MatematiskeVerktoyViTrenger/K2_algebra_beregning.ipynb}{\texttt{TRES0312 Python/Kap2\_MatematiskeVerktoyViTrenger/K2\_algebra\_beregning.ipynb}}.}

De vanlige regneoperasjonene er:

\begin{table}[H]
\centering
\begin{tabular}{|c|l|c|}
\hline
\rowcolor{blue!20!black!10}
\textbf{Symbol} & \textbf{Operasjon} & \textbf{Eksempel} \\
\hline
\texttt{+} & Addisjon       & \texttt{3 + 2 = 5}   \\
\texttt{-} & Subtraksjon    & \texttt{5 - 3 = 2}   \\
\texttt{*} & Multiplikasjon & \texttt{4 * 3 = 12}  \\
\texttt{/} & Divisjon       & \texttt{10 / 4 = 2.5}\\
\texttt{**}& Potens         & \texttt{2**3 = 8}    \\
\hline
\end{tabular}
\end{table}

\section{Skrive ut resultater: \texttt{print()}}

Funksjonen \texttt{print()} viser verdier på skjermen. Med en \emph{f-streng} kan vi blande tekst og variabler:

\begin{pythonkode}
E_k = 0.5 * m * v**2
print(f"Kinetisk energi: {E_k} J")
# Skriver ut:  Kinetisk energi: 25.0 J
\end{pythonkode}
\pyfilenote{Resten av samme kodecelle som i \href{\pyrepobase/Kap2_MatematiskeVerktoyViTrenger/K2_algebra_beregning.ipynb}{\texttt{TRES0312 Python/Kap2\_MatematiskeVerktoyViTrenger/K2\_algebra\_beregning.ipynb}}.}

Formatet \texttt{\{E\_k:.2f\}} begrenser til to desimaler. \texttt{\{theta:.1f\}} gir én desimal.

\section{NumPy — matematikk og arrays}

\textbf{NumPy} er et bibliotek for matematiske beregninger. Vi importerer det slik:

\begin{pythonkode}
import numpy as np
\end{pythonkode}
\pyfilenote{Denne importlinjen står øverst i de fleste notebook-filene i \href{https://github.com/kosmoben/TRES0312/tree/main/TRES0312\%20Python/Kap2_MatematiskeVerktoyViTrenger}{\texttt{TRES0312 Python/Kap2\_MatematiskeVerktoyViTrenger}}.}

Nyttige funksjoner:

\begin{table}[H]
\centering
\begin{tabular}{|l|l|}
\hline
\rowcolor{blue!20!black!10}
\textbf{Funksjon} & \textbf{Hva den gjør} \\
\hline
\texttt{np.sqrt(x)}       & Kvadratrot av $x$ \\
\texttt{np.sin(x)}        & Sinus ($x$ i radianer) \\
\texttt{np.cos(x)}        & Cosinus \\
\texttt{np.arctan2(y,x)}  & Vinkel i grader (bruk med \texttt{np.degrees}) \\
\texttt{np.degrees(x)}    & Gjør om fra radianer til grader \\
\texttt{np.linspace(a,b,n)}& $n$ jevnt fordelte tall mellom $a$ og $b$ \\
\texttt{np.array([1,2,3])} & Lager en array (liste med tall) \\
\texttt{np.dot(u, v)}     & Skalarprodukt \\
\texttt{np.linalg.norm(v)}& Lengden til en vektor \\
\hline
\end{tabular}
\end{table}

\section{Matplotlib — lage grafer}

\textbf{Matplotlib} brukes til å tegne grafer:

\begin{pythonkode}
import matplotlib.pyplot as plt

x = np.linspace(0, 10, 200)   # 200 punkter fra 0 til 10
y = x**2                       # y = x^2

plt.plot(x, y, color='royalblue', linewidth=2)
plt.xlabel('x')
plt.ylabel('y')
plt.title('Graf av y = x^2')
plt.grid(True, alpha=0.3)
plt.show()
\end{pythonkode}
\pyfilenote{Dette er et generelt eksempel og finnes ikke som egen fil i \texttt{TRES0312 Python} — prøv gjerne å skrive det inn selv i en ny notebook.}

Det er alt du trenger for å kjøre eksemplene i kapittel~1. Lykke til!

\end{document}
